% META: {"id":"TCOMPL-INEG-CO-004","chapitre":"TCOMPL-INEGALITES","type_objet":"corrige","exercice_ref":"TCOMPL-INEG-EX-004","capacites_codes":["C2"],"sources_inspiration":[],"status":"approved","origin":"generated","status_history":["generated","audited","accepted","approved"],"release_acceptance":"RELEASE_OWNER_BATCH_ACCEPTANCE","acceptance_closure_digest":"sha256:9b3ccf9a81c5520fbb7e03b7b2d3e7bf2057a02834b6d908bf3be5f49f3b553f","acceptance_receipt_digest":"sha256:4fad86f0282e0fa4e0419cca1ad6379ae7af5a280c04a4a90880f90a1c044242"}
% BEGIN-VERIFY
% from sympy import Rational
% assert Rational(1, 2)**4 < Rational(1, 2)**2
% END-VERIFY
\begin{corrige}{TCOMPL-INEG-EX-004}
\textbf{1.} Pour $0<x<1$, on a $x^4 < x^2$ (élever à une puissance plus grande un nombre entre $0$ et $1$ le rend plus petit). Donc $L_B(0{,}5) < L_A(0{,}5)$. Concrètement, $L_A(0{,}5)$ et $L_B(0{,}5)$ représentent la part de richesse détenue par les $50\%$ les plus pauvres, dans le pays A puis dans le pays B.

\textbf{2.} Comme $L_B(0{,}5) < L_A(0{,}5)$, les $50\%$ les plus pauvres détiennent une part de richesse plus faible dans le pays B : sa courbe de Lorenz est plus éloignée de la bissectrice, donc le pays B a la répartition la \textbf{plus inégalitaire}.
\end{corrige}
